\section{Opening Model}


This section organizes the central thesis of this paper — consciousness is not modeled as a generated object — as an opening model.

"Opening" here is not a metaphorical expression.

It is a structural description: when the fast layer, slow layer, plasticity gate, meaning gradient, and internal delay enter a certain relation, a log-referenceable passage is temporarily opened.

Consciousness is not a new object created inside.

Consciousness is the state in which referenceability has opened between already-running processing and already-sedimented logs.

\subsection{The Generation Assumption}


Most discussions of consciousness adopt a generation model, explicitly or implicitly.

In the generation model, consciousness is treated as something produced by some processing.

Under this view, the following questions naturally arise: which brain region generates consciousness; which circuit produces it; which computational quantity gives rise to it; how much consciousness is generated; can generated consciousness be stored or copied?

These questions seem natural at first glance.

But they all contain the presupposition of treating consciousness as a product: a generator, a generation process, and a generated object.

This three-term structure is the foundation of the generation model.

This paper does not adopt that presupposition.

Consciousness is not treated here as the product of processing; it is the state that opens when the relation between processing and logs enters a certain temporal condition.

\subsection{From Product to Relation}


In the generation model, consciousness is treated as a product.

In this paper, consciousness is treated as a relation.

That is, consciousness is the referencing relation established between fast predictive processing and slow log sedimentation.

For this relation to obtain, at least the following conditions are required:

\begin{enumerate}
\item Processing is running in the fast layer.
\item Referenceable logs have sedimented in the slow layer.
\item There is a finite speed difference $\Delta v$ between the two.
\item Coupling viscosity $\mu$ lies between complete synchrony and complete separation.
\item Internal delay $\tau$ is maintained as referenceable thickness.
\item Meaning gradient $\nabla\Phi$ is providing a referencing direction.
\item Effective torque-like quantity $\mathcal{T}_{\Phi}$ lies within the consciousness window.
\end{enumerate}

When these conditions are met, consciousness is not "produced."

Log referencing opens.

Consciousness is therefore not an output but an opening of a relation.

\subsection{Opening Through the Plasticity Gate}


The plasticity gate plays a central role in the opening model.

The plasticity gate is not a device that produces consciousness.

It is the boundary condition regulating whether fast-layer processing can sediment into the slow layer, and to what degree slow-layer logs influence current processing.

If the gate is completely closed, processing does not reach the slow layer. Even though processing runs, log referencing does not open.

If the gate is completely open, processing sediments excessively and tends to lose plasticity. Consciousness tends less to open freely and more to become fixed to specific logs or meaning wells.

What the opening model requires is not complete opening.

Rather, it is partial and temporary opening.

This opening is regulated by speed difference, coupling viscosity, internal delay, and meaning gradient:

\begin{equation}
\mathcal{T}_{\Phi} \sim \kappa \frac{\Delta v}{\mu} |\nabla \Phi|
\end{equation}

When $\mathcal{T}_{\Phi}$ lies within the consciousness window, the gate opens in a form appropriate for log referencing.

At that point, the conscious state obtains.

\subsection{Consciousness Opens, Qualia Appear}


As seen in the previous section, qualia are the live foreground appearing at the boundary where log referencing is in the process of sedimenting.

In the opening model, qualia too are not products.

Qualia are the phenomenal surface that appears at the boundary when consciousness opens.

For consciousness to open is for a referenceable passage to be established between the fast and slow layers.

At the boundary of that passage, processing has not yet completely become a log.

But it has not completely flowed away either.

This transitional state appears as qualia.

Qualia are therefore not the cause that produces consciousness.

Qualia are the live foreground observed at the boundary of the consciousness opening.

In the generation model, qualia invite the question "where are they produced?"

In the opening model, the question changes to: under what conditions does the boundary of log referencing appear as a live foreground?

This question can be described through speed difference, viscosity, delay, meaning gradient, and plasticity gate.

\subsection{Why Consciousness Cannot Be Stored}


From the perspective of the opening model, the idea of storing consciousness is structurally misplaced.

What is stored is not consciousness itself.

What is stored is the log structure updated while consciousness was open.

Equally, what is stored is not qualia themselves.

What is stored is the trace sedimented after qualia have passed through.

Consciousness obtains only while the state is open.

When the opening closes, the conscious state does not persist.

But its traces can remain in the slow layer.

This is why people can remember past experiences.

But what is remembered is not past consciousness itself.

It is the log that was transformed while past consciousness was open.

In this respect, the opening model distinguishes memory from consciousness.

Memory is stored. Consciousness is not stored.

Consciousness opens when stored logs are again referenced by current processing.

\subsection{Why Consciousness Cannot Be Copied}


For the same reason, consciousness cannot straightforwardly be copied.

What can be copied is structure, records, weights, logs, reports, representations.

But consciousness is not those objects themselves.

Consciousness is the state that opens when those objects enter a specific temporal relation.

Replicating the log structure at a given moment therefore does not guarantee that the replica has the same conscious state.

What is required is not merely log replication.

It is that fast processing, slow sedimentation, plasticity gate, speed difference, delay, meaning gradient, and non-closure loop operate in the same way.

Copying consciousness is therefore not data copying but reproduction of opening conditions.

In this respect, the opening model does not reduce consciousness to information content.

Consciousness is not content but the dynamic establishment of a relation.

\subsection{Opening and Non-Closure}


The opening model naturally corresponds to VED's non-closure principle.

In a completely closed structure, no opening exists.

In a completely open structure, sedimentation does not obtain.

What consciousness requires is the middle.

Tending toward closure but not closing completely. Open but not flowing away. Fixed but not over-fixed.

Within this range of temporal non-closure, log referencing opens.

The opening model therefore does not treat consciousness as an incompletely generated product.

It treats consciousness as the opening that obtains when non-closure has stabilized.

In this sense, rather than "consciousness is not generated," more precisely: "consciousness opens through not closing completely."

\subsection{Opening and Self}


The ego and self are also repositioned within this opening model.

The self is not the subject that owns consciousness.

The self is the referencing pattern that stabilizes as a result of log referencing being repeated and a certain terrain forming in the slow layer.

If consciousness is a temporary opening, the self is the terrain formed after that opening has repeatedly passed through.

The self is therefore not the cause of consciousness.

Rather, self-likeness sediments as consciousness repeatedly opens, logs are updated each time, and meaning gradients form.

In this view, the self is neither illusion nor absolute subject.

The self is the stable structure of the slow layer formed through the repetition of log referencing.

\subsection{Summary}


This section organized consciousness as an opening model.

\begin{itemize}
\item Consciousness is not a product.
\item Consciousness is the state in which referenceability has opened between fast processing and slow logs.
\item The plasticity gate is not a device that produces consciousness but a boundary condition regulating log-referencing opening.
\item Qualia are the live foreground appearing at the opening boundary.
\item What is stored is not consciousness but the logs updated while consciousness was open.
\item What can be copied is logs and structure, not consciousness itself.
\item Consciousness opens within the range of temporal non-closure, not at complete closure or complete opening.
\item The self is not the cause of consciousness but the slow-layer terrain formed through the repetition of opening.
\end{itemize}

The central thesis of this paper can therefore be restated as:

Consciousness is not modeled here as a generated object.

Consciousness obtains when log referencing opens within temporal non-closure.
